\begin{scanpage}{0191}{173}
correspond to the field radiated by the nonuniform part of the current. Comparing Expressions (27.04) and (27.03), we find that
\[
\left.\begin{aligned}\overline\Sigma_1^0&=\overline\Sigma^0,\quad\Sigma_1^0=0,\\
\overline\Sigma_2^0&=0,\quad\Sigma_2^0=-\overline\Sigma^0.
\end{aligned}\right\}\tag{27.07}
\]
Therefore, the field scattered in the direction towards the source ($\vartheta=\pi-\gamma$, $\varphi=-\pi/2$), will equal
\[
\left.\begin{aligned}E_x=-H_\vartheta&=\frac{ia}2[E_{0x}(\overline\Sigma^0+\overline\Sigma_1^1)+H_{0x}\Sigma_1^1]\frac{e^{ikR}}R,\\
E_\vartheta=H_x&=\frac{ia}2[E_{0x}\overline\Sigma_2^1-H_{0x}(\overline\Sigma^0-\Sigma_2^1)]\frac{e^{ikR}}R.
\end{aligned}\right\}\tag{27.08}
\]
This field's electric vector forms an angle $\beta$ with the yoz plane. The angle $\beta$ is determined by the equation
\[
\operatorname{tg}\beta=\frac{E_x}{E_\vartheta}=\frac{\overline\Sigma^0+\overline\Sigma_1^1-\Sigma_1^1\operatorname{ctg}\alpha}{\overline\Sigma^0-\Sigma_2^1+\overline\Sigma_2^1\operatorname{tg}\alpha}\operatorname{tg}\alpha.\tag{27.09}
\]
As a result, the desired angle $\delta$ which characterizes the depolarization magnitude will equal
\[
\delta=\alpha-\beta.\tag{27.10}
\]
Thus, the field from the nonuniform part of the current, separable ``in a pure form'' by means of a polarizer (\S~26), leads to depolarization of the scattered radiation.

Specific results from the depolarization calculation of waves reflected from certain bodies may be found, for example, in the works of Chytil [75--77] and Beckmann [78]. In particular, in Reference [77] it was shown that the depolarization effect on the effective scattering surface of convex bodies in practice may be neglected only with the condition $ka\gtrsim4$.
\end{scanpage}
